% 星群（综述）
% license CCBYSA3
% type Wiki

本文根据 CC-BY-SA 协议转载翻译自维基百科\href{https://en.wikipedia.org/wiki/Asterism_(astronomy)}{相关文章}。

\begin{figure}[ht]
\centering
\includegraphics[width=8cm]{./figures/c47b677da81e5ec4.png}
\caption{} \label{fig_Asteri_1}
\end{figure}
星群（Asterism）是指在天空中可见的恒星排列或星群图案。星群可以是任何被识别出的恒星形状或组合，因此它的概念比正式定义的 88 个星座更为广泛。星座是以星群为基础的，但与星群不同，星座是具有官方边界的明确区域，整体上覆盖了整个天空。

星群的形态从仅由几颗恒星构成的简单图案，到由大量恒星组成、横跨天空大片区域的复杂结构不等。构成星群的恒星可能是肉眼可见的明亮星，也可能是较暗、甚至需要望远镜才能观测到的星体，但它们通常在亮度上彼此相近。那些较大且较亮的星群对熟悉夜空的人而言具有重要的指引作用。

星群中恒星的排列并不一定反映它们之间的物理联系，而往往是由于观测视角造成的视觉效果。例如，“夏季大三角”仅仅是从地球观察到的一个表观组合，成员恒星之间在空间上并无关联；而“猎户座腰带”中的恒星则都属于猎户座 OB1 星协，“北斗七星”中的七颗星有五颗属于大熊座移动星群（Ursa Major Moving Group）。某些具有真实物理关联的星团，如“毕星团”（Hyades）或“昴星团”（Pleiades），既可以被视为独立的星群，也可能同时是更大星群的一部分。
\subsection{星群与星座的背景}
\begin{figure}[ht]
\centering
\includegraphics[width=8cm]{./figures/590b7d57cdce8a34.png}
\caption{} \label{fig_Asteri_2}
\end{figure}